\documentclass{beamer}
\usepackage{amsmath, amssymb}
\usepackage{lscape}
\usepackage{graphicx}

\begin{document}

%------------------------------------------------
\begin{frame}
\frametitle{Row Space and Column Space (Big Picture)}

Think of a matrix $A$ as a machine:
\[
A : \mathbb{R}^n \to \mathbb{R}^m
\]

Two natural questions:
\begin{itemize}
    \item What can its rows build?
    \item What can its columns build?
\end{itemize}

These give us:
\begin{itemize}
    \item Row Space
    \item Column Space
\end{itemize}

Both are just spans of vectors.
\end{frame}

%------------------------------------------------
\begin{frame}
\frametitle{Row Space}

If $A$ is $m \times n$, each row is a vector in $\mathbb{R}^n$.

Row Space = all linear combinations of the rows.

Example:
\[
A =
\begin{pmatrix}
1 & 2 & -3 \\
1 & 1 & 1
\end{pmatrix}
\]

Row space:
\[
\text{Span}\{(1,2,-3), (1,1,1)\}
\]

So it's all vectors of the form:
\[
\alpha(1,2,-3) + \beta(1,1,1)
\]

Geometrically:
A plane (or line) inside $\mathbb{R}^3$.
\end{frame}

%------------------------------------------------
\begin{frame}
\frametitle{Column Space}

Columns live in $\mathbb{R}^m$.

Column Space = all linear combinations of the columns.

Same matrix:
\[
A =
\begin{pmatrix}
1 & 2 & -3 \\
1 & 1 & 1
\end{pmatrix}
\]

Columns:
\[
\begin{pmatrix}1\\1\end{pmatrix},
\begin{pmatrix}2\\1\end{pmatrix},
\begin{pmatrix}-3\\1\end{pmatrix}
\]

Column space:
\[
\text{Span of those 3 vectors}
\]

Geometrically:
All outputs $Ax$ can ever produce.
\end{frame}

%------------------------------------------------
\begin{frame}
\frametitle{Very Important Interpretation}

\[
Ax = b
\]

This equation is solvable iff:

\[
b \in \text{Column Space of } A
\]

Translation:
The system is consistent exactly when $b$ is buildable
from the columns.

Columns = building blocks of outputs.
\end{frame}

%------------------------------------------------
\begin{frame}
\frametitle{Row Operations Do NOT Change Row Space}

If you row reduce a matrix:

\[
A \to U
\]

The row space stays the same.

Why?

Because row operations are just recombining rows.

So:
\begin{itemize}
    \item Reduce to RREF
    \item Nonzero rows form a basis
\end{itemize}
\end{frame}

%------------------------------------------------
\begin{frame}
\frametitle{Rank (Super Important)}

Rank = dimension of the row space.

How to compute it?

\begin{itemize}
    \item Row reduce
    \item Count nonzero rows
\end{itemize}

Example result:
If RREF has 2 nonzero rows,
\[
\text{rank}(A) = 2
\]

Rank = number of independent constraints.
Rank = number of pivots.
\end{frame}

%------------------------------------------------
\begin{frame}
\frametitle{Column Space and Rank}

Fact:
Row space and column space have the SAME dimension.

So:
\[
\dim(\text{Row Space}) = \dim(\text{Column Space})
\]

That number is the rank.

But careful:

Row reducing changes the column space!

So:
\begin{itemize}
    \item Use RREF to find pivot columns
    \item Go back to ORIGINAL matrix
    \item Those pivot columns form basis
\end{itemize}
\end{frame}

%------------------------------------------------
\begin{frame}
\frametitle{Null Space}

Null space = all solutions of:

\[
Ax = 0
\]

These are inputs that get crushed to zero.

Geometrically:
Directions that disappear under the transformation.

If there are free variables,
null space has dimension $> 0$.
\end{frame}

%------------------------------------------------
\begin{frame}
\frametitle{Rank-Nullity Theorem}

For $A$ with $n$ columns:

\[
\text{rank}(A) + \text{nullity}(A) = n
\]

Interpretation:

\begin{itemize}
    \item Rank = number of pivot variables
    \item Nullity = number of free variables
    \item Total variables = pivots + free
\end{itemize}

It’s just counting.
\end{frame}

%------------------------------------------------
\begin{frame}
\frametitle{When Is $Ax=b$ Unique?}

If columns are:
\begin{itemize}
    \item Independent → at most one solution
    \item Span $\mathbb{R}^m$ → solution exists for all $b$
\end{itemize}

If square matrix:

Unique solution for every $b$
iff
\[
A \text{ is invertible}
\]

iff
\[
\text{rank}(A) = n
\]
\end{frame}

%------------------------------------------------
\begin{frame}
\frametitle{Big Geometric Picture}

Matrix as transformation:

\[
A : \mathbb{R}^n \to \mathbb{R}^m
\]

\begin{itemize}
    \item Column space = output space
    \item Null space = inputs that collapse to zero
    \item Rank = how many dimensions survive
    \item Nullity = how many dimensions collapse
\end{itemize}

Rank + Nullity = total input dimensions.
\end{frame}

%------------------------------------------------
\begin{frame}
\frametitle{Final Intuition}

Think in CS terms:

\begin{itemize}
    \item Columns = features
    \item Rank = number of independent features
    \item Null space = redundant directions
    \item Row reduction = simplifying constraints
\end{itemize}

Row space = constraints.
Column space = outputs.
Null space = invisible inputs.
Rank = surviving dimensions.
\end{frame}

%------------------------------------------------
\end{document}